
\section{Proposal}
\label{Pro}
\subsection{Overview}

In this section, we present the overall flow of our proposal.
%, depicted in Fig.~\ref{fig-flow}. The blue boxes represent precomputation, while tan boxes represent runtime.

Library precomputation:
\begin{itemize}
\item \textbf{(Generate 3bit phase polynomials)} : Precompute library of phase polynomials
  of Boolean functions of three bits or fewer using the Boolean Fourier transform.
\item \textbf{(Synthesize 3bit phase subcircuits)} : The phase polynomials are then
  sequenced using matroid partitioning from~\cite{bib-amy-matroid} and synthesized as
  subcircuits. The results are stored in hard disk, to be loaded into memory at runtime.
\end{itemize}

Synthesis:
\begin{itemize}
\item \textbf{(Decompose $f$ into ESOP expression)} : At runtime, a Boolean function is
  read in and converted to an ESOP expression.
\item \textbf{(Minimize the number of $n \geq 4$ cubes in ESOP)} : The ESOP is optimized
  using methods such as EXORCISM-4~\cite{bib-exorcism} to reduce the number of cubes of
  more than 4 variables. Our extensions to such methods is explained in
  \secref{Pro:Minimize}.
\item \textbf{(Map the $n \leq 3$ cubes to the 3 bit library)} : Takes all cubes of
  $n \leq 3$ and maps them to appropriate functions from the precomputed 3 bit library, concatenating the
  corresponding subcircuit with the phase oracle so far.
\item \textbf{(Synthesize $n \geq 4$ cubes)} : Synthesizes all cubes of $n \geq 4$ using the
  method described in Sec.~\ref{Pro:n4}.
\item \textbf{(Combine $n \geq 4$ and $n \leq 3$)} : $n \geq 4$ is concatenated with
  $n \leq 3$ in series.
\item \textbf{(Optimize for T-count and T-depth)} : The combined circuit is optimized for T-count and T-depth using
  any algorithm that optimizes T-depth in CNOT+T circuits~\cite{bib-amy-matroid}. 
\end{itemize}
%% \begin{figure}[t]
%%   \centering
%%   \scalebox{1.0} {
%%     \input{img/flowchart.tex}
%%   }
%%   \caption{The overall flow of the proposed method}
%%   \label{fig-flow}
%% \end{figure}

\subsection{Minimize the number of $n \geq 4$ cubes in ESOP}
\label{Pro:Minimize}

To minimize the number of $n \geq 4$ cubes in the ESOP, the proposal adapts EXORCISM-4~\cite{bib-exorcism}.
EXORCISM-4 is a heuristic method of minimizing the number of cubes in a given ESOP expression. Extended
discussion of its details is deferred to the cited paper. 

We modify it for our purposes by changing
the cost function \texttt{CountCubesInExactPseudoKro} that EXORCISM-4 uses for its heuristic, which uses a plain sum
of cubes to assess the cost of the arrangement that its heuristic encounters. We introduce a cost function
that instead uses a weighted sum of the cubes, where each cube is weighted by its T-count (i.e. a 2 literal cube
has weight 0) for $n \leq 3$ and $2^n$ for $n \geq 4$ (its T-count when synthesized according to the method in
Sec.~\ref{Pro:n4}). This penalizes terms of $n \geq 4$, while terms of $n < 3$ are free.


\subsection{Synthesize $n \geq 4$ cubes}
\label{Pro:n4}

To synthesize cubes of $n \geq 4$ variables, we utilize the construction in Fig.~\ref{fig-phase-n4}. Here,
$U_F$ and $U_F^{\dagger}$ are $F$-controlled-NOT gates, meaning that they apply X to the target bit when
the Boolean function $F$ is true. We add an ancilla bit and use the relative-phase MCT synthesis
method from~\cite{bib-rtof_maslov} to enable the implementation of such $F$-controlled-NOT gates with a
reduced number of T-gates. We defer detailed discussion of that method to the cited paper. This resulting
phase oracle has T-count $2^{n}$, and T-depth $2n+1$. Note that this means that larger cubes can
dominate the T-depth of a given circuit, so it is imperative that the minimum of them be found for
any given Boolean function.

%% \subsection{Optimize for T-count and T-depth}
%% \label{Pro:Tpar}

%% While this paper's techniques can be applied to any method that optimizes using CNOT+T circuits, we opt to target
%% the T-par algorithm from~\cite{bib-amy-matroid}. This is a variant of the matroid partitioning algorithm that deals
%% with optimization of Clifford+T circuits by heuristically separating the circuit into CNOT+T partitions along the
%% H-gates in the circuit. This deals with the circuits synthesized in Sec.~\ref{Pro:n4} as well as the CNOT+T circuits
%% the $n \leq 3$ processing step produces. We defer more detailed discussion to the cited paper.



